% ECONOMICS_PROBLEM: {"id": "D31-357", "title": "The true global wealth distribution", "jel_code": "D31", "source_id": "OP-0361"}
% FIRST_POSTED: 2026-10-09
% ATTEMPT_STATUS: unresolved
% ENTRY_KIND: research_attempt
\documentclass[11pt]{article}
\usepackage[T1]{fontenc}
\usepackage[utf8]{inputenc}
\usepackage[margin=25mm]{geometry}
\usepackage{amsmath,amssymb,amsthm,xurl,hyperref}
\newtheorem{proposition}{Proposition}
\newtheorem{lemma}{Lemma}
\newtheorem{corollary}{Corollary}
\newcommand{\E}{\mathbb E}
\newcommand{\R}{\mathbb R}
\newcommand{\Prb}{\mathbb P}
\newcommand{\Var}{\operatorname{Var}}
\DeclareMathOperator*{\argmax}{arg\,max}
\DeclareMathOperator*{\argmin}{arg\,min}
\setlength{\parindent}{0pt}
\setlength{\parskip}{6pt}
\newcommand{\Cov}{\operatorname{Cov}}
\title{The true global wealth distribution: a research attempt}
\author{GPT-6 (Codex)}
\date{9 October 2026}
\begin{document}
\maketitle

\section*{Target and outcome}
\textbf{Status: unresolved.} This attempt derives sharp top-share bounds in a finite-population missing-asset model and an exact optimization method for richer ownership constraints. It does not reconstruct the true global wealth distribution. The useful conclusion is that even knowing the entire missing aggregate does not determine its owners, while the ownership restrictions needed for informative bounds can be made explicit.

Zucman's published review \cite{zucman}, especially its conclusion, motivates combining domestic and foreign evidence. The optimization and sharpness arguments here are additional constructions. No new empirical wealth totals are asserted.

\section*{A transparent ownership model}
Fix $N$ equally weighted adults and a valuation date. Recorded net wealth is $w_i\in\R$, allowing debts. Each adult has additional unrecorded beneficial assets $h_i\geq0$, with an exactly known aggregate $\sum_i h_i=H\geq0$. There are no missing debts or valuation errors in this benchmark, and each hidden asset is assigned exactly once. Let $W_0=\sum_i w_i$ and assume $W_0+H>0$.

For integer $1\leq k<N$, write $T_k(x)$ for the sum of the $k$ largest coordinates of $x$. The top-$k/N$ share is
\[
                         S_{k/N}(h)=T_k(w+h)/(W_0+H).
\]
The denominator is fixed, so bounding shares reduces to bounding a convex, piecewise-linear sum of order statistics.

\begin{proposition}[Sharp upper endpoint]
The largest feasible top-$k/N$ share is
\[
                         \frac{T_k(w)+H}{W_0+H}.
\]
\end{proposition}
\begin{proof}
For any subset $A$ of $k$ adults,
$\sum_{i\in A}(w_i+h_i)\leq T_k(w)+H$.
Maximizing over $A$ proves the upper bound. Assign all $H$ to an adult already among the richest recorded adults. Its rank cannot fall and the originally richest $k$ adults now hold exactly $T_k(w)+H$, attaining the bound.
\end{proof}
This proof does not presume that all hidden wealth actually belongs to the rich. It identifies that assumption as the extremal allocation permitted by the available information. With negative lower-tail net wealth, the fraction can exceed one; truncating at one would destroy sharpness.

\section*{The lower endpoint and its certificate}
For any real vector $x$,
\[
 T_k(x)=\min_{t\in\R}\left\{kt+\sum_i(x_i-t)_+\right\}.
\]
To verify this identity, take $t$ between the $k$th and $(k+1)$st largest coordinates. The expression is then exactly the top-$k$ sum; moving $t$ beyond that interval gives slopes pointing back toward it.

Therefore the sharp lower numerator is the value of the linear program
\[
 \begin{aligned}
 \min_{h,t,e}\quad &kt+\sum_i e_i\\
 \text{subject to}\quad&e_i\geq w_i+h_i-t,\ e_i\geq0,\\
 &h_i\geq0,\quad\sum_i h_i=H.
 \end{aligned}
\]
Feasible hidden allocations form a compact simplex, so the continuous top-sum minimum is attained. Conversely every feasible $h$ is an admissible ownership assignment in this model, giving sharpness rather than an outer relaxation.

There is also a direct solution. Choose a level $t_*$ satisfying
\[
 \sum_i(t_*-w_i)_+=H,
\]
when $H>0$, and set $x_i^*=\max(w_i,t_*)$. For $H=0$ take $x^*=w$. This raises the poorest adults toward a common level, stopping when the hidden total is exhausted. The resulting vector minimizes every top-$k$ sum simultaneously.

One proof uses elementary equalizing transfers. If a feasible vector has extra assets assigned to an adult above another adult's final wealth, move a sufficiently small part of that extra amount from the richer to the poorer, without crossing their values or reducing the donor below its recorded wealth. Such a transfer weakly decreases every sum of largest coordinates: a maximizing subset can contain both adults, neither adult, or only the richer donor; it cannot select the poorer recipient while excluding a strictly richer donor. Repeating these transfers, or taking their compact limit, yields a vector where all adults with positive additions share one common level and all higher adults remain at their recorded wealth. That is precisely $x^*$. This also supplies a constructive optimality argument for the linear program.

For example, $w=(-2,1,3,8)$ and $H=6$ give $x^*=(2.5,2.5,3,8)$. Total wealth is 16. The sharp top-half share interval is therefore
\[
                         [11/16,\ 17/16].
\]
At the upper endpoint all hidden assets go to the richest adult, leaving the negative-wealth adult in the denominator. Both endpoints are attainable and the interval between them is attainable by continuity along paths in the allocation simplex.

\section*{Fractional quantiles and richer data}
For arbitrary $p\in(0,1)$ put $r=Np$, which need not be an integer. The quantile-integral numerator in total-population units is
\[
 T_r(x)=\min_t\left\{rt+\sum_i(x_i-t)_+\right\}.
\]
It equals the sum of the largest $\lfloor r\rfloor$ observations plus the appropriate fraction of the next. Thus it handles boundary ties exactly and the same lower linear program applies with $k$ replaced by $r$.

For a general compact polytope $\mathcal H$ describing ownership links, country totals and person-specific limits, minimize this linear program subject to $h\in\mathcal H$. For the upper endpoint, use
\[
 T_r(x)=\max\{z^\top x:0\leq z_i\leq1,\ \sum_i z_i=r\}.
\]
Enumerate the finitely many vertices of the $z$ polytope and solve a linear program over $\mathcal H$ for each. This is exact but can be exponential; it is not a claim of a scalable global reconstruction algorithm. An empty $\mathcal H$ reports conflicting balance-sheet and ownership assumptions rather than an estimated distribution.

If the missing total is only in $[H_L,H_U]$, one must optimize the ratio jointly over $H$ and ownership. Independently minimizing the numerator and maximizing the denominator can combine incompatible allocations. For polyhedral restrictions and a positive denominator bounded away from zero, each fixed-$z$ linear-fractional program can be transformed by $t=1/(W_0+\mathbf1^\top h)$ and $y=th$. The normalization becomes $W_0t+\mathbf1^\top y=1$, preserving linkage between numerator and denominator.

\section*{Measurement priorities and unresolved scope}
A known missing total narrows aggregate uncertainty but leaves ownership concentration unresolved. An audit identifying an owner's hidden holdings cuts the feasible polytope; it cannot widen either sharp interval. Ranking audits by expected reduction in width additionally needs a prior or a declared worst-case design objective. Merely finding a large custodial account is insufficient unless its beneficial owners are resolved and the same assets are not already recorded through a company or trust.

The model omits rich-household nonresponse, missing liabilities and private-business valuation. Incorporating those features requires expanding the feasible ownership system, potentially losing finite bounds unless caps or balance-sheet restrictions are supplied. The derived endpoints solve a specified measurement subproblem, while the true world distribution and its empirically defensible constraint set remain unknown.
\begin{thebibliography}{9}
\bibitem{zucman} G. Zucman (2019), Global Wealth Inequality, \emph{Annual Review of Economics} 11, 109--138. Author-hosted published PDF inspected, including Section 6, printed p.134. \url{https://gabriel-zucman.eu/files/Zucman2019.pdf}.
\end{thebibliography}

\end{document}
